\documentclass[12pt,oneside]{article}
\usepackage[T1]{fontenc}
\usepackage[utf8]{inputenc}

\usepackage{geometry}
\geometry{letterpaper, top=1.25in, bottom=1.25in, left=1.25in, right=1.25in}
\usepackage{amsmath, amsthm, amssymb, mathtools}
\usepackage[colorlinks=true,linkcolor=black,citecolor=black,urlcolor=black,
  pdftitle={Error as Obstruction},pdfauthor={Flyxion}]{hyperref}
\usepackage{setspace}
\onehalfspacing
\usepackage{parskip}
\setlength{\parskip}{0.5\baselineskip}
\usepackage{titlesec}
\titleformat{\section}{\large\bfseries}{\thesection.}{0.6em}{}
\titleformat{\subsection}{\normalsize\bfseries}{\thesubsection.}{0.6em}{}
\usepackage{enumitem}
\setlist[itemize]{noitemsep,topsep=2pt}

\theoremstyle{plain}
\newtheorem{theorem}{Theorem}[section]
\newtheorem{lemma}[theorem]{Lemma}
\newtheorem{corollary}[theorem]{Corollary}
\newtheorem{proposition}[theorem]{Proposition}
\theoremstyle{definition}
\newtheorem{definition}[theorem]{Definition}
\theoremstyle{remark}
\newtheorem{remark}[theorem]{Remark}

\newcommand{\X}{\mathcal{X}}
\newcommand{\M}{\mathcal{M}}
\newcommand{\Sh}{\mathcal{F}}
\newcommand{\D}{D}
\newcommand{\Hcheck}{\check{H}}
\newcommand{\Sig}{\Sigma}
\newcommand{\rhotau}{\rho_\tau}
\newcommand{\Asoft}{\mathcal{A}_{\mathrm{soft}}}
\newcommand{\HessSig}{\mathrm{Hess}_{\Sigma}}

\begin{document}

\begin{center}
  {\LARGE\bfseries Error as Obstruction}\\[1.2em]
  {\large Flyxion}\\[0.4em]
  {\normalsize Independent Researcher}\\[1.8em]
  {\small\today}
\end{center}
\bigskip

\begin{abstract}
\noindent
We develop a unified geometric account of error, learning, and event-sourced
computation in which failure of coherence is identified not with deviation from a
ground-truth value but with a topological impossibility of extension. The central
claim is that an event log, a training dataset, or any preserved history
generates a constraint sheaf over a domain; admissible states are precisely
global sections of that sheaf; and what is ordinarily called an error is a
cohomological obstruction to the existence of such a section. We formalise this
claim through the \emph{Error as Obstruction} theorem, which equates the
convergence of stochastic projected relaxation dynamics with the triviality of a
derived cohomology class. The framework synthesises the RSVP field-space
construction, the Spherepop event-calculus,
and the TARTAN sheaf-consistency architecture into a single
programme: intelligence is the organisation of an overcomplete constraint manifold
under boundary preservation, and error is the failure of that manifold to admit
a globally coherent section.
\end{abstract}

\newpage
\tableofcontents
\newpage

%=============================================================
\section{Introduction}
%=============================================================

\subsection{The Bottleneck Assumption and Its Discontents}

For most of its history, the theory of learning has been organised around
scarcity. The core assumption is that the capacity to represent is limited, that
experience exceeds what can be retained, and that generalisation is therefore the
art of productive forgetting---of compressing many instances into a compact rule.
Under this view, memory and generalisation stand in fundamental tension: a system
that preserves too much detail overfits, mistaking the texture of particular
examples for the structure of the underlying process. The standard prescription
follows directly: reduce, simplify, discard.

This assumption is so deeply embedded that it shapes not only formal models but
practical pedagogy. A struggling student is told to take fewer courses. A
craftsperson is advised to specialise. A programmer is counselled to keep the
codebase small. The recommendation in every case is subtraction: reduce the
inputs until the system stabilises.

The present work begins from a different direction. There is a class of learning
systems---and a class of learners---for which this prescription is not merely
unhelpful but actively wrong. For these systems, adding more structured
experience does not increase the burden; it increases the coherence. Consider a
builder who knows electrical, plumbing, construction, finishing, and installation.
When tasked with sorting materials, organising offcuts, or evaluating whether a
scrap fitting will be useful, that person is not performing five separate
assessments. Every fragment is already indexed by function, hazard, sequence, and
reuse potential. The narrow specialist sees disconnected objects; the
cross-trained worker sees a structured landscape. A narrow assistant either asks
constantly or proceeds blindly---not because they lack facts in isolation, but
because they lack the surrounding geometry that would tell them what matters,
what can wait, what is hazardous, and what they do not yet know they do not know.

The same pattern appears at different scales. Switching between subjects when
studying is not avoidance; it is coordinate-frame rotation, a way of keeping the
system active while preventing any single track from becoming locally stale. In
a namespace already crowded with thousands of projects, one does not need to
remember every prior name; enough exposure to the distribution internalises it,
so that a new choice can be evaluated as a position in a landscape whose rough
topology is already known. The choice is no longer arbitrary but navigational.

These examples share a common structure. In each case, the system is operating
above a threshold beyond which additional structured input reshapes the space of
possible organisation rather than merely occupying it. Memorisation and
generalisation are not in tension because the system is not a fixed container
filling up; it is a geometry being refined. Capacity here is not the size of a
buffer but the dimensionality of an internal coordinate system, and that
dimensionality can \emph{increase} with exposure.

\subsection{The Formal Problem}

The standard picture of learning identifies error with a scalar distance between
prediction and target. This scalar is differentiable, and gradient descent drives
it toward zero. The picture is clean, but it conceals the structural question
that motivates the present work: \emph{why does perfect interpolation not destroy
generalisation?} In the excess-capacity regime---neural networks with far more
parameters than training examples, event stores with richer constraint
vocabularies than the log they hold---the system memorises everything and yet
continues to extrapolate sensibly~\cite{Zhang2017,Belkin2019,DubovaSloman2026}.
Compression cannot explain this; the system has not compressed.

The answer we develop here is geometric rather than statistical. Once the system
reaches exact fit, learning does not stop; it continues as motion \emph{within}
the admissible set, driven by stochastic relaxation dynamics that seek regions of
low curvature. In this regime, error is not a gradient to be cancelled but a
structural impossibility: the log or dataset has imposed locally consistent
constraints that cannot be glued into any globally coherent configuration. That
impossibility is an obstruction in the sheaf-theoretic sense, and it is the
central object of the paper.

The novelty of the claim is not that memorisation and generalisation can coexist.
Any practitioner who has worked in a richly cross-constrained domain already
knows this from the inside. The novelty is that their coexistence is not
accidental. It is a \emph{predictable consequence} of operating in a regime where
the space of possible representations is larger than the data that constrains
it---a regime with a precise geometric characterisation, a well-defined dynamical
theory, and a natural notion of failure.

That notion of failure is the obstruction of the title. When a history cannot be
made globally coherent---when locally consistent constraints cannot be glued into
any smooth global configuration---the system does not merely give a wrong answer.
It inhabits a space in which no coherent answer exists. Error, on this account,
is not a mismeasurement. It is a topological fact about the structure of the
record.

The exposition proceeds as follows. Section~\ref{sec:setting} establishes the
mathematical setting: the RSVP field space, the admissible fit set, and the
three geometric regimes of learning. Section~\ref{sec:dynamics} introduces the
soft action functional and derives projected gradient flow on regular strata.
Section~\ref{sec:stochastic} extends the construction to stochastic dynamics and
defines the stationary density of wisdom. Section~\ref{sec:entropy} develops the
role of tangential entropy and shows that selection in the excess regime is a
volumetric phenomenon. Section~\ref{sec:sovereignty} articulates the principle of
Log Sovereignty. Section~\ref{sec:boundary} formalises the log as a boundary
operator and proves historical closure. Section~\ref{sec:spherepop} identifies
the Spherepop event-calculus as a discrete realisation of these principles.
Section~\ref{sec:singular} analyses singular strata and degenerate geometry as
the precursor to full obstruction. Section~\ref{sec:sheaf} introduces the
constraint sheaf and reformulates admissibility as global section existence.
Section~\ref{sec:obstruction} states and proves the main theorem together with
the Stability--Obstruction Dichotomy. Section~\ref{sec:diagnostic} isolates the
triple diagnostic signature. Section~\ref{sec:classification} classifies
obstruction types and develops a hierarchy of repair operations.
Section~\ref{sec:viscoelastic} makes the viscoelastic picture exact.
Section~\ref{sec:ml} connects the framework to contemporary machine learning.
Section~\ref{sec:conclusion} states the central thesis in its strongest form.
An appendix recounts the phenomenological motivation in more concrete terms for
readers whose primary interest is practical rather than formal.

%=============================================================
\section{Mathematical Setting}
\label{sec:setting}
%=============================================================

\subsection{The RSVP Field Space}

Let $\Omega \subset \mathbb{R}^n$ be a domain. A \emph{field configuration} is
a triple $\Phi = (\phi, \mathbf{v}, S)$ where $\phi : \Omega \to \mathbb{R}$ is
a scalar density field, $\mathbf{v} : \Omega \to \mathbb{R}^n$ is a vector flow
field, and $S : \Omega \to \mathbb{R}_{\geq 0}$ is an entropy density.
The configuration space $\X$ is a Fr\'echet manifold whose local charts are
determined by Sobolev norms on each component.

A readout operator $R : \X \to \mathbb{R}^m$ extracts observable predictions
from a configuration. The experience set $\D = \{(x_i, y_i)\}_{i=1}^N$
specifies $N$ input-output pairs. We write $\hat{y}_i = R(\Phi)(x_i)$ for the
prediction at $x_i$.

\begin{definition}[Admissible Fit Set]
The \emph{admissible fit set} for $\D$ is
\[
  \M(\D) = \bigl\{\,\Phi \in \X \;\big|\; R(\Phi)(x_i) = y_i
  \;\text{for all}\; i = 1,\dotsc,N\,\bigr\}.
\]
\end{definition}

$\M(\D)$ is the zero-level set of the constraint map $\Phi \mapsto
(R(\Phi)(x_i) - y_i)_{i=1}^N$ and, for smooth $R$, is generically a
submanifold of $\X$ whose codimension equals the number of active constraints.

\subsection{The Three Geometric Regimes}

\noindent\textbf{Constrained regime.} The constraints over-determine $\X$:
$\M(\D) = \emptyset$, and the system can only minimise a soft penalty.

\medskip\noindent\textbf{Sufficient regime.} The number of constraints matches
the effective degrees of freedom: $\M(\D)$ is generically a discrete set of
isolated points.

\medskip\noindent\textbf{Excess regime.} The system has far more degrees of
freedom than constraints: $\M(\D)$ is a high-dimensional manifold, and learning
continues as navigation \emph{within} $\M(\D)$. This is not a pathology but a
phase transition in the geometry of the problem~\cite{Belkin2019,DubovaSloman2026}.

%=============================================================
\section{Projected Gradient Flow}
\label{sec:dynamics}
%=============================================================

\subsection{The Soft Action Functional}

Let $\Asoft : \X \to \mathbb{R}_{\geq 0}$ be a smooth, coercive \emph{soft
action functional} encoding inductive preferences beyond exact fit.

\begin{definition}[Restricted Hessian]
For $\Phi \in \M(\D)$, let $T_\Phi\M(\D)$ denote the tangent space and let
$\Pi_\Phi$ denote orthogonal projection onto it. The \emph{restricted Hessian}
is
\[
  \HessSig\Asoft(\Phi) = \Pi_\Phi\,D^2\Asoft(\Phi)\,\Pi_\Phi.
\]
\end{definition}

\begin{definition}[Flat Region]
A stratum $U \subset \M(\D)$ is \emph{flat} if $\HessSig\Asoft(\Phi) = 0$ for
all $\Phi \in U$, and \emph{low-curvature} if its operator norm is bounded
by $\varepsilon \ll 1$ uniformly on $U$.
\end{definition}

\subsection{Projection-Relaxation Dynamics}

The intramanifold dynamics are governed by the \emph{projected gradient flow}:
\[
  \dot{\Phi}(t) = -\,\Pi_{\Phi(t)}\,\nabla\Asoft(\Phi(t)).
\]
This suppresses motion normal to $\M(\D)$---directions that would violate log
constraints---while allowing free motion in tangential directions. The normal
component of the gradient is the \emph{restriction bias}; the tangential
component is the \emph{soft inductive bias}.

\begin{proposition}[Monotonicity]
\label{prop:mono}
Along any trajectory of the projected gradient flow,
$\tfrac{d}{dt}\Asoft(\Phi(t)) \leq 0$, with equality if and only if $\Phi(t)$
is a critical point of $\Asoft|_{\M(\D)}$.
\end{proposition}

\begin{proof}
By the chain rule,
$\tfrac{d}{dt}\Asoft(\Phi) = \langle\nabla\Asoft(\Phi),\dot\Phi\rangle
= -\|\Pi_\Phi\nabla\Asoft(\Phi)\|^2 \leq 0$.
\end{proof}

%=============================================================
\section{Stochastic Relaxation and the Stationary Density of Wisdom}
\label{sec:stochastic}
%=============================================================

\subsection{The Stochastic Projected Relaxation Equation}

We consider the \emph{stochastic projected relaxation equation}
\[
  d\Phi_t = -\Pi_{\Phi_t}\,\nabla\Asoft(\Phi_t)\,dt
  + \sqrt{2\tau}\;\Pi_{\Phi_t}\,dW_t,
\]
where $W_t$ is a cylindrical Wiener process on $\X$ and $\tau > 0$ is a
temperature parameter.

\subsection{The Fokker--Planck Equation on a Smooth Stratum}

For a compact smooth stratum $\Sig \subset \M(\D)$, the density satisfies
\[
  \partial_t\rhotau
  = \mathrm{div}_\Sig(\rhotau\,\nabla_\Sig\Asoft)
  + \tau\,\Delta_\Sig\rhotau,
\]
where $\nabla_\Sig$ and $\Delta_\Sig$ are the intrinsic gradient and
Laplace--Beltrami operator on $\Sig$. As $\tau \to 0$, the stationary measure
concentrates on the sub-level sets of $\Asoft|_\Sig$~\cite{Otto2001,Friston2010}.

\begin{definition}[Wisdom]
\label{def:wisdom}
A configuration $\Phi \in \M(\D)$ is \emph{wise} if it lies in a low-curvature
region $U \subset \M(\D)$ on which $\rhotau$ concentrates uniformly for all
sufficiently small $\tau > 0$. A system is wise if its long-run state
distribution is concentrated in flat strata.
\end{definition}

The geometric principle is: \emph{a stable world-state is not one that is simply
true, but one that is truest across the widest range of possible tangential
interpretations.}

%=============================================================
\section{Entropy, Measure, and Selection}
\label{sec:entropy}
%=============================================================

\subsection{Entropy as Tangential Volume}

The stochastic projected relaxation dynamics induce a probability measure
$\rhotau$ supported on $\M(\D)$. While $\Asoft$ governs energy, the geometry of
$\M(\D)$ itself determines entropy. In the excess regime the admissible set
carries a nontrivial volume form, and the system's long-run behaviour reflects a
balance between energetic preference and geometric multiplicity.

\begin{definition}[Tangential Entropy]
Let $\Sig \subset \M(\D)$ be a smooth stratum. The \emph{tangential entropy}
at $\Phi \in \Sig$ is
\[
  \mathcal{S}_{\mathrm{tan}}(\Phi)
  = -\int_{U_\Phi} \rhotau(\Psi)\,\log\rhotau(\Psi)\,d\Psi,
\]
where $U_\Phi$ is a neighbourhood in $\Sig$ and $d\Psi$ is the Riemannian
volume element induced by the ambient metric on $\X$.
\end{definition}

Flat regions maximise $\mathcal{S}_{\mathrm{tan}}$ by supporting large volumes
of near-equivalent configurations. Stochastic dynamics therefore do not merely
descend energy gradients; they select regions of maximal admissible volume.

\subsection{Selection as Measure Concentration}

The stationary density
\[
  \rhotau(\Phi) \;\propto\; \exp\!\Bigl(-\frac{\Asoft(\Phi)}{\tau}\Bigr)
\]
assigns weight not only by energy but by the measure of neighbourhoods. Regions
of low curvature dominate because they occupy disproportionately large volume in
$\Sig$. This establishes the following principle.

\begin{proposition}[Volumetric Selection]
\label{prop:volume}
In the excess regime, the stationary measure $\rhotau$ concentrates on
low-curvature strata of $\M(\D)$ as $\tau \to 0$, independently of whether those
strata contain the global minimum of $\Asoft$.
\end{proposition}

\begin{proof}
By Laplace's method on the manifold $\Sig$, the mass assigned to a region
$U \subset \Sig$ under $\rhotau$ scales as
$\exp(-\Asoft(\Phi_0)/\tau) \cdot \mathrm{vol}(U) \cdot (1 + O(\tau))$,
where $\Phi_0$ is the energy minimiser in $U$ and $\mathrm{vol}(U)$ is the
Riemannian volume. A flat region with large volume can dominate a sharp minimum
with small volume whenever the volume ratio exceeds
$\exp((\Asoft(\Phi_{\mathrm{sharp}}) - \Asoft(\Phi_{\mathrm{flat}}))/\tau)$.
\end{proof}

Generalisation is therefore a volumetric phenomenon. A configuration is stable
not because it minimises $\Asoft$ absolutely, but because it lies in a region
that remains invariant under perturbation~\cite{Hochreiter1997,Keskar2017}. The
excess-capacity regime is exactly the regime in which such regions exist and are
accessible; the constrained regime collapses them to isolated points.

%=============================================================
\section{Log Sovereignty}
\label{sec:sovereignty}
%=============================================================

\begin{definition}[Log Sovereignty]
A log $\D$ exercises \emph{sovereignty} over a system if $\D$ is treated as a
boundary condition rather than a compressible object: $\M(\D)$ is defined by
exact equality constraints derived from $\D$, and the system's dynamics remain
within $\M(\D)$ at all times.
\end{definition}

Under Log Sovereignty, forgetting is not a mechanism of generalisation. The log
is preserved entirely---a choice with roots in the thermodynamic principle that
erasure of information is the only irreversible operation~\cite{Landauer1961,Bennett1982}.
What generalises is the system's capacity to maintain coherence across the
high-dimensional admissible structure: to find and remain in regions where
$\Asoft$ is low. This dissolves the standard tension between memorisation and
generalisation---in the excess regime, $\M(\D)$ is large enough to contain both,
and the stochastic relaxation selects the stable configurations by navigating
within the constraint manifold toward its flattest regions.

%=============================================================
\section{Boundary Conditions and Historical Closure}
\label{sec:boundary}
%=============================================================

\subsection{The Log as a Boundary Operator}

The event log $\D$ is not merely a set of constraints but a \emph{boundary
operator} acting on the field space $\X$. It defines the admissible set as a
boundary-constrained subspace:
\[
  \M(\D) = \ker\partial_\D,
\]
where $\partial_\D : \X \to \mathbb{R}^N$ maps configurations to constraint
violations, $\partial_\D(\Phi) = (R(\Phi)(x_i) - y_i)_{i=1}^N$.

This perspective makes explicit that $\D$ does not describe the state; it
\emph{restricts} the space of possible states. The state is never the log; the
log is the boundary within which the state lives.

\subsection{Historical Closure}

\begin{definition}[Historical Closure]
A system exhibits \emph{historical closure} if all admissible configurations
are determined by $\D$ as a boundary condition, and all dynamics preserve
membership in $\M(\D)$ for all time.
\end{definition}

Under historical closure, the past is not stored as data to be retrieved but
enforced as geometry. The admissible manifold is the closure of history under
the dynamics.

\begin{proposition}[Invariance under Projected Relaxation]
\label{prop:invariance}
Historical closure is equivalent to the invariance of $\M(\D)$ under the
projected relaxation dynamics.
\end{proposition}

\begin{proof}
The projected gradient flow $\dot\Phi = -\Pi_\Phi\nabla\Asoft(\Phi)$ removes
all components of the gradient normal to $\M(\D)$ by construction. Therefore,
if $\Phi(0) \in \M(\D)$, then $\Phi(t) \in \M(\D)$ for all $t \geq 0$.
Conversely, any dynamics that preserve $\M(\D)$ must remove normal components,
which is exactly the projection condition.
\end{proof}

This formalises Log Sovereignty at the level of dynamics rather than intent.
The log is not a policy that the system chooses to respect; it is a geometric
constraint that the dynamics cannot exit. Sovereignty is structural, not
volitional.

\subsection{Remark on Discrete Logs}

In the continuous field-theoretic setting, $\partial_\D$ is a smooth operator
and $\M(\D)$ is a smooth submanifold at regular values. In the discrete
event-sourced setting of Section~\ref{sec:spherepop}, the same structure holds
locally: each event $e_i$ contributes one constraint, and the admissible set is
the intersection of the corresponding zero-level sets. The boundary operator
formalism is therefore not an artefact of the continuous setting but the common
structural core of both.

%=============================================================
\section{Spherepop as Discrete Projected Relaxation}
\label{sec:spherepop}
%=============================================================

\subsection{The Event Log as Boundary Condition}

In the Spherepop calculus, a world-state is defined by its \emph{identity as
event history}. Let $\D = (e_1, \dotsc, e_N)$ be an ordered event sequence.
Each $e_i$ imposes the constraint that $\Phi$ is consistent with $e_i$. The
admissible world-state set
\[
  \M(\D) = \{\,\Phi \in \X \mid \Phi\text{ is consistent with }e_i
  \text{ for all }i\,\}
\]
is the construction of Section~\ref{sec:setting}. Reconstruction of state from
the log is not deterministic replay but a search for an admissible configuration.

\subsection{The Dual-Rate Projection-Relaxation Cycle}

When a new event $e_{N+1}$ arrives, the admissible set transitions from
$\M(\D)$ to $\M(\D \cup \{e_{N+1}\})$. This update has two distinct phases.

\medskip\noindent\textbf{The Normal Snap (fast).} The system projects its
current configuration onto the updated admissible set:
\[
  \Phi' = \operatorname*{argmin}_{\Psi \in \M(\D\cup\{e_{N+1}\})}
  \|\Phi - \Psi\|_\X.
\]
This enforces logical consistency with the new event immediately and makes the
system \emph{truthful} with respect to the log.

\medskip\noindent\textbf{The Tangential Drift (slow).} Starting from $\Phi'$,
the system runs stochastic projected relaxation on $\M(\D\cup\{e_{N+1}\})$,
seeking low-curvature regions. This makes the system \emph{stable} under
perturbation.

The Normal Snap alone does not determine behaviour. A configuration can be
admissible while residing in a high-curvature, brittle region of $\M(\D)$. The
Tangential Drift converts correctness into wisdom. Systems that treat event
replay as deterministic and terminal---rather than as initialisation for ongoing
relaxation---conflate the two phases and thereby sacrifice stability.

%=============================================================
\section{Singular Strata and Degenerate Geometry}
\label{sec:singular}
%=============================================================

\subsection{Stratification of the Admissible Set}

In general, $\M(\D)$ is not a smooth manifold but a \emph{stratified space}:
\[
  \M(\D) = \bigsqcup_k \Sig_k,
\]
where each $\Sig_k$ is a smooth manifold of dimension $k$. Singularities arise
at configurations where the rank of the constraint map $\partial_\D$ drops below
its generic value: these are the points at which multiple constraint surfaces
are tangent, or at which the Jacobian of $\partial_\D$ loses rank. The regular
stratum is the open dense subset on which the rank is maximal; all other strata
are lower-dimensional and singular.

\subsection{Breakdown of Tangential Structure}

At a singular point $\Phi \in \M(\D)$, the tangent space $T_\Phi\M(\D)$ is not
well-defined as a linear subspace. Instead one has a \emph{tangent cone}: a
closed set of possible limiting tangent directions that need not span a
subspace. The orthogonal projection $\Pi_\Phi$ becomes multivalued or
discontinuous, and the projected gradient flow fails to be well-posed.

\begin{definition}[Geometric Degeneracy]
A configuration $\Phi \in \M(\D)$ is \emph{degenerate} if the projection
operator $\Pi_\Phi$ fails to define a unique, continuous map from $T_\Phi\X$
into any subspace.
\end{definition}

Degeneracy is the geometric precursor to obstruction. It is the point at which
the system can no longer distinguish admissible tangential directions from
inadmissible normal ones; the two have become mixed.

\subsection{Dynamical Consequences at Singular Strata}

Near singular strata, stochastic projected relaxation exhibits characteristic
pathological behaviour:

\begin{itemize}
\item rapid fluctuations between competing tangent directions as the cone
      has multiple extreme rays;
\item amplification of noise from the loss of curvature control, since the
      restricted Hessian $\HessSig\Asoft(\Phi)$ is ill-defined and cannot
      provide a restoring force;
\item failure of local linearisation, so that first-order approximations to
      the dynamics are qualitatively incorrect.
\end{itemize}

These effects constitute an early warning signal: singular strata are the
regions in which the dynamical consequences of potential obstruction first
manifest, even before the full cohomological obstruction described in
Section~\ref{sec:obstruction} is reached.

\begin{remark}
Regular strata support smooth relaxation; singular strata disrupt it. The
transition between them is the geometric onset of error. A system whose
trajectories repeatedly enter singular strata without settling is exhibiting
the precursor signature of an obstructed log.
\end{remark}

\subsection{Relation to Obstruction}

The singular strata of $\M(\D)$ are precisely the configurations at which the
constraint sheaf (Section~\ref{sec:sheaf}) is most fragile. Near a singularity,
the local sections of $\Sh_\D$ over nearby patches become difficult to glue,
because the local geometry no longer provides a consistent normal direction to
guide the gluing. Full obstruction---the non-triviality of the \v{C}ech
class---is the global version of the local degeneracy diagnosed here.

%=============================================================
\section{Constraint Sheaves and Global Sections}
\label{sec:sheaf}
%=============================================================

\subsection{The Constraint Sheaf}

Let $\mathcal{U} = \{U_\alpha\}$ be a cover of the index set $\{1,\dotsc,N\}$
by overlapping patches, where each $U_\alpha$ represents a local context in
which the events $\{e_i : i \in U_\alpha\}$ are simultaneously relevant.

\begin{definition}[Constraint Sheaf]
The \emph{constraint sheaf} $\Sh_\D$ assigns to each patch $U_\alpha$ the
locally admissible configurations:
\[
  \Sh_\D(U_\alpha)
  = \bigl\{\,\Phi|_{U_\alpha} \in \X|_{U_\alpha} \;\big|\;
  \Phi|_{U_\alpha}\text{ is consistent with }e_i
  \text{ for all }i \in U_\alpha\,\bigr\},
\]
together with the natural restriction maps on overlaps.
\end{definition}

\begin{proposition}
\label{prop:section}
$\M(\D)$ is nonempty if and only if $\Sh_\D$ admits a global section.
\end{proposition}

\subsection{The \v{C}ech Cohomology of Incoherence}

The obstruction to extending locally consistent sections to a global section is
measured by $\Hcheck^1(\mathcal{U};\Sh_\D)$. A class
$[\sigma] \in \Hcheck^1(\mathcal{U};\Sh_\D)$ is represented by a 1-cocycle
$\sigma = (\sigma_{\alpha\beta})$ satisfying the cocycle condition on triple
overlaps but not cobounding---not writable as $\sigma_{\alpha\beta} =
\tau_\alpha - \tau_\beta$ for any assignment of local sections $\tau_\alpha$.

When $[\sigma] \neq 0$, local sections cannot be glued into a global one. A
logical error, a hallucination, or a coherence failure is, under this analysis, a
non-trivial \v{C}ech class: a mismatch on overlaps that cannot be absorbed by
any choice of global configuration~\cite{BottTu1982,GelfandManin1996,AbramskyBrandenburger2011,Ghrist2014}.

%=============================================================
\section{The Error as Obstruction Theorem}
\label{sec:obstruction}
%=============================================================

\begin{theorem}[Error as Obstruction]
\label{thm:main}
Let $\D$ be an event log inducing a constraint sheaf $\Sh_\D$ over a cover
$\mathcal{U}$, and let $\M(\D)$ be the admissible fit set of global sections of
$\Sh_\D$. Let $\Asoft$ be a smooth, coercive soft action functional on $\X$, and
consider the stochastic projected relaxation dynamics on regular strata of
$\M(\D)$.

Then the following are equivalent:
\begin{enumerate}[label=(\roman*)]
  \item There exists a nonempty regular stratum $\Sig \subset \M(\D)$ supporting
        stable stochastic projected relaxation with a well-defined stationary
        measure.
  \item The constraint sheaf $\Sh_\D$ admits a global section, equivalently
        $[\sigma] = 0$ in $\Hcheck^1(\mathcal{U};\Sh_\D)$.
\end{enumerate}

Conversely, if $[\sigma] \neq 0$ in $\Hcheck^1(\mathcal{U};\Sh_\D)$, then:
\begin{enumerate}[label=(\alph*)]
  \item $\M(\D)$ is either empty or consists entirely of singular configurations
        at which $\Pi_{T_\Phi\M(\D)}$ is not well-defined.
  \item The stochastic projected relaxation fails to converge to any stable
        distribution; trajectories exhibit persistent high-action behaviour.
  \item The action functional satisfies
        $\displaystyle\inf_{\Phi \in \X}\,\Asoft(\Phi) > 0$.
\end{enumerate}
\end{theorem}

\begin{proof}
\noindent\emph{(i) $\Rightarrow$ (ii).}
A regular stratum is a smooth submanifold; its existence implies $\M(\D)
\neq\emptyset$. By Proposition~\ref{prop:section}, this implies $[\sigma] = 0$.

\medskip\noindent\emph{(ii) $\Rightarrow$ (i).}
Suppose $[\sigma] = 0$, so a global section $\Phi^* \in \M(\D)$ exists. The
implicit function theorem, applied to the constraint map
$\Phi \mapsto (R(\Phi)(x_i) - y_i)$, shows that near any regular value the
admissible set is a smooth submanifold. This submanifold is nonempty and
compact in appropriate Sobolev topologies when $\Asoft$ is coercive. Standard
results on stochastic differential equations on compact manifolds guarantee the
existence of a unique stationary measure for the stochastic projected relaxation.

\medskip\noindent\emph{Converse.}
Suppose $[\sigma] \neq 0$. Then $\M(\D)$ contains no smooth point: any candidate
global section fails the compatibility condition on at least one overlap,
introducing a geometric defect. At such a defect, $T_\Phi\M(\D)$ does not exist
and $\Pi_{T_\Phi\M(\D)}$ is undefined. The stochastic trajectories, unable to
project onto a consistent tangent, accumulate action rather than dissipating
it---the persistent high-action behaviour of~(b). The lower bound in~(c) follows
from the coercivity of $\Asoft$ and the absence of admissible configurations at
which $\Asoft$ can reach zero.
\end{proof}

\begin{remark}
Theorem~\ref{thm:main} identifies error with the nonexistence of a smooth global
section compatible with the log. This is not approximation error, prediction
error, or deviation from a target value; it is structural impossibility. The
system does not fail to find the right answer; it fails to inhabit a space in
which any answer is coherent.
\end{remark}

\begin{theorem}[Stability--Obstruction Dichotomy]
\label{thm:dichotomy}
Let $\D$ be an event log inducing a constraint sheaf $\Sh_\D$, and let $\Asoft$
be a smooth, coercive soft action functional on $\X$. Consider the stochastic
projected relaxation dynamics restricted to $\M(\D)$.

Then exactly one of the following holds:
\begin{enumerate}[label=(\roman*)]
  \item \textbf{Stability.} The cohomology class vanishes, $[\sigma] = 0$ in
        $\Hcheck^1(\mathcal{U};\Sh_\D)$. There exists a nonempty regular stratum
        $\Sig \subset \M(\D)$ on which the dynamics admit a unique stationary
        measure $\rhotau$, and trajectories converge in distribution toward
        low-curvature regions of $\Sig$.
  \item \textbf{Obstruction.} The cohomology class is non-trivial,
        $[\sigma] \neq 0$. The admissible set $\M(\D)$ contains no smooth
        stratum supporting stable dynamics, and stochastic projected relaxation
        fails to converge, exhibiting persistent action floor, tangent
        instability, and non-localising trajectories.
\end{enumerate}
These cases are mutually exclusive and exhaustive.
\end{theorem}

\begin{proof}
\emph{Mutual exclusivity.} By Theorem~\ref{thm:main}, existence of a regular
stratum with stable dynamics is equivalent to $[\sigma] = 0$, while
non-triviality of the cohomology class implies the absence of such strata.
The two conditions therefore cannot hold simultaneously.

\emph{Exhaustivity.} For any $\D$, either $\Sh_\D$ admits a global section or
it does not. In the first case $\M(\D)$ contains smooth points and supports
projected relaxation with a stationary measure by standard results on stochastic
flows on compact manifolds~\cite{Hairer2009}. In the second case $\M(\D)$ is empty or entirely
singular; the projection operator fails to define a consistent tangential flow,
producing the dynamical signature of case~(ii). No third case exists, since any
intermediate behaviour would require smooth strata incompatible with the
global section obstruction.
\end{proof}

\begin{corollary}[Wisdom--Error Separation]
\label{cor:wisdom_error}
Wisdom and error are separated by a topological invariant. A system is capable
of sustaining a stable notion of wisdom if and only if $[\sigma] = 0$. When
$[\sigma] \neq 0$, no refinement of the relaxation dynamics alone can produce
stability; modification of the log is necessary.
\end{corollary}

\begin{remark}
The dichotomy eliminates the possibility of systems that are ``almost coherent''
in any stable sense. Apparent intermediate cases correspond to long-lived
transient behaviour near singular strata (Section~\ref{sec:singular}), not to a
third stable regime. The system either admits a globally coherent organisation or
it does not; the topological invariant $[\sigma]$ decides which.
\end{remark}

%=============================================================
\section{The Triple Diagnostic Signature}
\label{sec:diagnostic}
%=============================================================

The proof of Theorem~\ref{thm:main} shows that obstruction is registered not by
a single scalar but by a \emph{triple signature}.

\medskip\noindent\textbf{1.~Persistent action floor.} In the unobstructed case,
$\Asoft(\Phi_t) \to 0$ as $t \to \infty$. In the obstructed case,
$\inf_{\Phi\in\X}\Asoft(\Phi) > 0$: the action does not decay but oscillates
above a floor set by the cohomological defect.

\medskip\noindent\textbf{2.~Tangent instability.} At any candidate configuration
near the defect, the restricted Hessian $\HessSig\Asoft(\Phi)$ is
ill-conditioned: the projection operator cannot produce a consistent tangential
direction because the manifold has a topological tear.

\medskip\noindent\textbf{3.~Non-convergence of stochastic trajectories.} In the
unobstructed case, the Fokker--Planck density on $\Sig$ converges to the Gibbs
measure $\rhotau \propto \exp(-\Asoft/\tau)$. In the obstructed case, no such
measure exists; trajectories wander without localising.

These signals are individually observable and jointly diagnostic. A system
exhibiting all three is not merely incorrect but incoherent in the topological
sense, and no local adjustment to its configuration can resolve the failure.

\begin{corollary}[Locality of Repair]
\label{cor:repair}
If $[\sigma]$ is supported on a single patch $U_\alpha$, modification of
$\{e_i : i \in U_\alpha\}$ is sufficient to restore admissibility. If the
support is global, the log requires structural revision.
\end{corollary}

%=============================================================
\section{Classification of Obstruction Types}
\label{sec:classification}
%=============================================================

\noindent\textbf{Local obstructions.} A 1-cocycle $\sigma$ is local if it is
cohomologous to a cocycle supported on a single overlap $U_\alpha \cap U_\beta$.
It arises from a pair of events whose combined constraints are mutually
incompatible and can be resolved by replacing one event or refining the cover.
This is the computational analogue of a merge conflict.

\medskip\noindent\textbf{Global obstructions.} The cohomology class is
non-trivial but not representable by a cocycle on any single overlap. The log is
globally inconsistent even though every pair of events is pairwise compatible.
Resolution requires adding new events that thread the inconsistency, or revising
the cover structure.

\medskip\noindent\textbf{Persistent obstructions.} The class is non-trivial in
$\Hcheck^k(\mathcal{U};\Sh_\D)$ for all $k \geq 1$: no refinement of the cover
can make the log globally consistent. Such logs must be partially retracted.

\subsection{A Hierarchy of Repair Operations}

The classification yields a hierarchy of repair operations ordered by increasing
cost. \emph{Event substitution} replaces a single offending event and resolves
local obstructions. \emph{Cover refinement} splits a patch to separate
conflicting constraints, potentially trivialising a local cocycle. \emph{Log
extension} adds new events whose constraints thread a global inconsistency.
\emph{Log retraction} removes a connected family of events, eliminating a
persistent obstruction at the cost of erasing part of the history---under Log
Sovereignty, the last resort.

%=============================================================
\section{Viscoelastic Logic}
\label{sec:viscoelastic}
%=============================================================

The tangent space decomposes as
\[
  T_\Phi\X = N_\Phi\M(\D) \oplus T_\Phi\M(\D).
\]
The normal component $N_\Phi\M(\D)$ is elastically suppressed, with stiffness
$\lambda \gg 1$. The tangential component $T_\Phi\M(\D)$ is viscously damped by
the stochastic relaxation, with viscosity coefficient $\tau$.

In the unobstructed case, the system behaves like a viscoelastic medium: stress
introduced by new events redistributes and dissipates. The elastic component
snaps the system to the constraint manifold; the viscous component drifts it
toward its most stable configuration.

In the obstructed case, the medium has a topological tear. Stress cannot be
distributed because the manifold does not support a continuous deformation from
the current configuration to any stable one. The elastic component deposits the
system near the defect; the viscous component cannot carry it away. This is
\emph{viscoelastic failure}: the material tears rather than deforming.

Viscoelastic logic is therefore not a metaphor but the exact dynamical regime
characterised by Theorem~\ref{thm:main}. The triple signature of
obstruction---action floor, tangent instability, trajectory non-convergence---is
the mechanical signature of a torn medium.

%=============================================================
\section{Connection to Machine Learning}
\label{sec:ml}
%=============================================================

\noindent\textbf{Double descent.} The double descent curve corresponds to the
transition through the three geometric regimes. The constrained regime is the
standard bias--variance setting~\cite{Belkin2019}. At the interpolation
threshold, $\M(\D)$ is near-discrete and highly curved, producing the variance
peak. In the excess regime, $\M(\D)$ is a smooth high-dimensional manifold, and
relaxation toward flat regions produces the second descent~\cite{Zhang2017}.

\medskip\noindent\textbf{Flat minima and implicit regularisation.} Stochastic
gradient descent with small batch sizes finds flatter minima because flat regions
have larger volume and greater measure under the Gibbs distribution; the
stochastic dynamics naturally select them~\cite{Hochreiter1997,Keskar2017,Amari1998}.

\medskip\noindent\textbf{Benign overfitting.} A model that achieves zero
training loss while generalising well navigates within $\M(\D)$ in a
low-curvature region~\cite{DubovaSloman2026}. Small perturbations to the input
do not move the system out of the flat region.

\medskip\noindent\textbf{Hallucination.} A hallucination is a configuration that
appears admissible under the fast Normal Snap but resides near a cohomological
defect. The system generates locally consistent outputs that cannot be glued into
a globally coherent section~\cite{AbramskyBrandenburger2011}. The diagnostic is
not the content of the output but the geometry of the neighbourhood: high
restricted Hessian, persistent action floor, stochastic non-convergence.

In each case, the machine learning phenomenon is a special instance of
intramanifold dynamics under stochastic projected relaxation. The framework
provides a geometric language in which existing empirical results become theorems.

%=============================================================
\section{Conclusion}
\label{sec:conclusion}
%=============================================================

The unification of Spherepop, RSVP,
and TARTAN developed here is a single continuous programme:
a redefinition of learning, computation, and state as geometric processes
constrained by preserved history.

\medskip\noindent\textbf{Central thesis.} \emph{Intelligence is not the
reduction of information but the maintenance and coherent organisation of it.
The capacity to do so depends not on limiting representation but on navigating
the geometry of excess. Error is not a deviation from ground truth but a
topological impossibility: the failure of a history to admit a globally coherent
configuration.}

\medskip
The thesis has four structural components. \emph{Truth} is the existence of a
global section of the constraint sheaf: constraints glued together without
remainder. \emph{Intelligence} is the capacity to navigate within that section
under the projected gradient flow. \emph{Wisdom} is the occupancy of
low-curvature strata---configurations stable under tangential perturbation.
\emph{Error} is the non-triviality of the \v{C}ech class: the log's constraints
cannot be made globally coherent by any choice of configuration. The appropriate
response to error is not correction---adjustment of a scalar---but
\emph{repair}: identification and resolution of the cohomological defect that
prevents the log from sustaining a coherent world-state.

The title is therefore not rhetorical. Error is obstruction: the precise
mathematical object that prevents the construction of a globally coherent global
section from a locally consistent event log. And the geometry of that
obstruction---its locality or globality, its persistence under cover
refinement---is what determines whether the log can be repaired, and how.

\bigskip
\noindent\textbf{Acknowledgements.} This work is developed within the ongoing
programme on RSVP, Spherepop, and TARTAN.

%=============================================================
\newpage
\appendix
\section{Phenomenological Motivation: Learning Without a Ceiling}
\label{app:phenom}
%=============================================================

This appendix offers a more concrete account of the intuitions behind the paper
for readers whose primary interest is in learning practice rather than in the
formal machinery. It is not part of the proof structure. It is the record that
gave rise to the question.

\subsection*{The Capacity That Grows}

The standard prescription for a struggling learner is reduction. Take fewer
courses. Narrow the focus. Let the overloaded buffer drain. This prescription
rests on a model: the mind as a container of fixed size, filling up as experience
accumulates, degrading when overfull. Under that model, relief requires
subtraction.

The model is incomplete. A student who routinely takes more courses than
required---and finds that performance improves rather than decays---is not a
counterexample to some statistical regularity. They are evidence that a different
regime is possible: one in which the constraint set is large enough that
additional experience does not compete with existing structure but
\emph{extends} it. The brain, under those conditions, is not a container. It is
a geometry. Courses do not fill it; they refine its coordinate system.

The same effect appears in skilled manual work. A builder who knows electrical,
plumbing, structural framing, finishing, and site management does not merely know
more facts than a single-trade specialist. The cross-domain knowledge creates a
dense field of relevance in which every object is already situated. A piece of
scrap wire is not just metal; it is indexed by gauge, application, code
implications, reuse potential, and whether it belongs to a live circuit. The
narrow specialist, lacking that surrounding structure, must either ask or
proceed blindly. They do not even know what they do not know, and so they make
errors that seem absurd to the practitioner who can see the whole landscape.

The student who switches subjects when bored is doing something structurally
similar. Each domain provides a coordinate system for pattern recognition. When
one domain is locally exhausted---when fatigue or repetition has suppressed
gradient---another domain provides fresh structure. The subject-switching is not
distraction; it is manifold traversal. The learner is not disorganised. They are
maintaining a coupled system of representations, each of which stabilises the
others.

\subsection*{Naming as Navigation}

A concrete version of the same principle appears in how one assigns names to
projects in a large and active shared namespace. Someone who has observed
thousands of repositories does not remember them individually. But they have
internalised the distribution: which names are saturated (every variation of
\texttt{core}, \texttt{browser}, \texttt{parser}, \texttt{base} is taken or
nearly so), which are empty, which are ambiguous, and which carry connotations
that will mislead. The act of naming a project becomes navigational rather than
arbitrary. When a collision occurs---when a new project appears using the same
name---the experienced practitioner moves their own project rather than insisting
on territory. The name is not a stake; it is a coordinate. What matters is
clarity in the shared space.

This is what it looks like to operate within a richly populated field. The
information retained is not a list of entries but a topology. The system holds a
model of the landscape, and each new observation refines that model rather than
merely adding to a count.

\subsection*{Cross-Domain Structure and the Invariants It Produces}

What repeated traversal of many domains produces, over time, is not merely
familiarity with many things. It produces \emph{invariants}---patterns that
recur across contexts, recognisable because they have been encountered in enough
variations that their structure is visible through the surface differences. This
is the mechanism behind what practitioners sometimes call intuition or
judgment: not a mysterious extra-cognitive faculty, but the accumulated result of
having encountered enough structured variation that the relevant invariants are
quickly activated by new instances.

The formal account in the body of this paper treats learning as motion on an
admissible manifold toward regions of low curvature. The phenomenological
account says the same thing in experiential language: the practitioner with
extensive cross-domain exposure is not searching for an answer from scratch. They
are already near a flat region of the space, because their prior experience has
already organised the surrounding geometry. New inputs do not require heavy
computation to integrate; they slot into an already-structured landscape.

\subsection*{The Limit Case: When Accumulation Outpaces Integration}

The one genuine risk in this regime is not accumulation per se but the
possibility that the rate of accumulation outpaces the rate of integration. If
inputs arrive faster than the system can form connections among them, they
accumulate as isolated tokens rather than as nodes in a connected field. That is
the only scenario in which the bottleneck model's prescription becomes relevant:
when the system is, in fact, operating in a disconnected way.

The diagnostic is not the count of inputs but the connectivity of the
representations they generate. A system in which many inputs are forming
connections---each new project illuminating previously opaque aspects of earlier
ones---is not overloaded. A system in which inputs accumulate without forming
connections is. The intervention in the second case is not necessarily
subtraction; it may be restructuring: slowing down long enough to allow
connections to form, or deliberately rotating between domains to expose latent
commonalities that were not visible under a single lens.

The paper's formal language for this is the constraint sheaf and its sections.
Locally consistent constraints that cannot be glued globally are the formal image
of knowledge that has not integrated. Persistent action floor, tangent
instability, and trajectory non-convergence are the signatures of a system that
is not yet coherent---not overloaded, but not yet organised. The remedy, in each
case, is not less experience but more structure in how experience is held.

\newpage
\begin{thebibliography}{99}

\bibitem{DubovaSloman2026}
M.~Dubova and S.~J.~Sloman,
\newblock ``Excess Capacity Learning,''
\newblock \emph{Behavioral and Brain Sciences}, 2026, pp.~1--77.
\newblock doi:10.1017/S0140525X2610510X.

\bibitem{Belkin2019}
M.~Belkin, D.~Hsu, S.~Ma, and S.~Mandal,
\newblock ``Reconciling modern machine-learning practice and the classical
  bias--variance trade-off,''
\newblock \emph{Proceedings of the National Academy of Sciences}, vol.~116,
  no.~32, pp.~15849--15854, 2019.

\bibitem{Zhang2017}
C.~Zhang, S.~Bengio, M.~Hardt, B.~Recht, and O.~Vinyals,
\newblock ``Understanding deep learning requires rethinking generalization,''
\newblock \emph{International Conference on Learning Representations}, 2017.

\bibitem{Hochreiter1997}
S.~Hochreiter and J.~Schmidhuber,
\newblock ``Flat minima,''
\newblock \emph{Neural Computation}, vol.~9, no.~1, pp.~1--42, 1997.

\bibitem{Keskar2017}
N.~S.~Keskar, D.~Mudigere, J.~Nocedal, M.~Smelyanskiy, and P.~T.~P.~Tang,
\newblock ``On large-batch training for deep learning: Generalization gap and
  sharp minima,''
\newblock \emph{International Conference on Learning Representations}, 2017.

\bibitem{MacLane1998}
S.~Mac~Lane,
\newblock \emph{Categories for the Working Mathematician},
\newblock Springer, 2nd edition, 1998.

\bibitem{BaezStay2011}
J.~C.~Baez and M.~Stay,
\newblock ``Physics, topology, logic and computation: A Rosetta Stone,''
\newblock in \emph{New Structures for Physics}, Springer, 2011, pp.~95--172.

\bibitem{Lurie2009}
J.~Lurie,
\newblock \emph{Higher Topos Theory},
\newblock Princeton University Press, 2009.

\bibitem{BottTu1982}
R.~Bott and L.~W.~Tu,
\newblock \emph{Differential Forms in Algebraic Topology},
\newblock Springer, 1982.

\bibitem{GelfandManin1996}
S.~I.~Gelfand and Y.~I.~Manin,
\newblock \emph{Methods of Homological Algebra},
\newblock Springer, 1996.

\bibitem{Spivak2014}
D.~I.~Spivak,
\newblock \emph{Category Theory for the Sciences},
\newblock MIT Press, 2014.

\bibitem{Jacobson1995}
T.~Jacobson,
\newblock ``Thermodynamics of spacetime: The Einstein equation of state,''
\newblock \emph{Physical Review Letters}, vol.~75, no.~7, pp.~1260--1263,
  1995.

\bibitem{Verlinde2011}
E.~Verlinde,
\newblock ``On the origin of gravity and the laws of Newton,''
\newblock \emph{Journal of High Energy Physics}, vol.~2011, no.~4, 2011.

\bibitem{Landauer1961}
R.~Landauer,
\newblock ``Irreversibility and heat generation in the computing process,''
\newblock \emph{IBM Journal of Research and Development}, vol.~5, no.~3,
  pp.~183--191, 1961.

\bibitem{Bennett1982}
C.~H.~Bennett,
\newblock ``The thermodynamics of computation---a review,''
\newblock \emph{International Journal of Theoretical Physics}, vol.~21,
  pp.~905--940, 1982.

\bibitem{Lloyd2000}
S.~Lloyd,
\newblock ``Ultimate physical limits to computation,''
\newblock \emph{Nature}, vol.~406, pp.~1047--1054, 2000.

\bibitem{Friston2010}
K.~Friston,
\newblock ``The free-energy principle: A unified brain theory?,''
\newblock \emph{Nature Reviews Neuroscience}, vol.~11, no.~2, pp.~127--138,
  2010.

\bibitem{Amari1998}
S.~Amari,
\newblock ``Natural gradient works efficiently in learning,''
\newblock \emph{Neural Computation}, vol.~10, no.~2, pp.~251--276, 1998.

\bibitem{Villani2009}
C.~Villani,
\newblock \emph{Optimal Transport: Old and New},
\newblock Springer, 2009.

\bibitem{Otto2001}
F.~Otto,
\newblock ``The geometry of dissipative evolution equations: The porous medium
  equation,''
\newblock \emph{Communications in Partial Differential Equations}, vol.~26,
  pp.~101--174, 2001.

\bibitem{Hairer2009}
M.~Hairer,
\newblock ``Introduction to stochastic partial differential equations,''
\newblock \emph{arXiv preprint arXiv:0907.4178}, 2009.

\bibitem{AbramskyBrandenburger2011}
S.~Abramsky and A.~Brandenburger,
\newblock ``The sheaf-theoretic structure of non-locality and contextuality,''
\newblock \emph{New Journal of Physics}, vol.~13, 2011.

\bibitem{Ghrist2014}
R.~Ghrist,
\newblock \emph{Elementary Applied Topology},
\newblock Createspace, 2014.

\end{thebibliography}

\end{document}
